%================================================================
\chapter{Introduction}
\label{ch:introduction}
%================================================================

\section{Motivation}
\label{sec:motivation}

Water distribution networks deliver drinking water to hundreds of
millions of people daily. The pipes and pumps are only part of it:
a modern network is a Cyber-Physical System (CPS), watched over by hundreds of
pressure and flow sensors that feed into a Supervisory Control and Data Acquisition (SCADA) system.
Operators use those readings to balance demand, catch pipe bursts,
and schedule maintenance. Wrong readings produce wrong responses.

The same connectivity that makes SCADA systems useful is
what an attacker needs. Industrial control systems were
historically air-gapped; today they are routinely connected to
corporate networks and the internet for remote access and centralised
monitoring. The Stuxnet worm in 2010~\cite{langner2011stuxnet} showed
that an attacker who reaches a control system can cause physical
damage without ever touching hardware. Water infrastructure has not
escaped attention: the Oldsmar, Florida incident in February
2021~\cite{oldsmar2021} saw an attacker remotely elevate the sodium
hydroxide concentration in a water treatment plant to 111 times the
safe level. The attack was caught by a vigilant operator who happened
to be watching the screen. Automated detection would have been faster
and more reliable.

Detecting anomalies in WDN sensor data is not simply a
matter of checking whether a reading falls outside a fixed range.
Pressure at a junction is correlated with pressure at every other
junction through the physics of the pipe network; a reading that is
unusual at one sensor may be perfectly consistent with a simultaneous
change at a neighbouring sensor, or it may be a falsification
designed to look that way. Classical detection methods based on
residuals from a hydraulic model require an accurate model and full
observability, neither of which holds in practice. Machine learning
methods that ignore the graph structure discard exactly the relational
information that distinguishes a real anomaly from a correlated
fluctuation.

Graph neural networks are a natural fit for this problem. A \wdn{}
is a graph: junctions are nodes, pipes are edges. A \gnn{} propagates
information along that graph through message passing, so each node's
embedding automatically reflects its neighbourhood hydraulics. The
question is not whether \gnn{}s are a reasonable choice but how to
make them robust against attackers who know the model and try to
fool it.

\section{Problem Statement}
\label{sec:problem}

This thesis addresses the following problem:

\begin{quote}
  \textit{Given a water distribution network monitored by pressure
  and flow sensors, where some sensors may be missing and some may be
  falsified by an adversary, train a model that simultaneously
  reconstructs missing readings and detects the falsified sensors,
  and remains reliable when the adversary adapts to the model.}
\end{quote}

The problem has two coupled sub-tasks. Reconstruction asks the model
to predict the true pressure and flow values at sensors whose
readings are absent. Anomaly detection asks it to output a binary
flag at each sensor indicating whether the current reading has been
manipulated. The two tasks share a backbone because reconstruction
quality determines how well the model understands the network state,
and that understanding underpins anomaly detection.

Robustness to adversarial adaptation is the hard part. A model
trained on a fixed set of hand-crafted attacks fits a distribution;
an attacker who knows that distribution can craft perturbations that
stay outside it. Self-play is one answer: train an attacker and a
defender jointly so that the defender's training distribution is
never fixed.

\section{Research Contributions}
\label{sec:contributions}

There are four contributions, and they are best read as a chain:
each one is the answer to a problem the previous one could not
solve.

\textbf{A spatial GNN baseline.} A GraphSAGE backbone with four
heads (pressure and flow reconstruction, pressure and flow anomaly)
brings pressure reconstruction MAE down by more than an order of
magnitude against classical state estimation (pseudo-inverse and
weighted least squares) on Net1. The point of this contribution is
mostly negative: it works well except on replay, and that failure
sets up the next step.

\textbf{A temporal extension that handles replay.} Wrapping each
snapshot in a GRU over a six-step window gives the model history.
Six steps is not arbitrary; it is about the least history that lets
the model notice a reading has stopped moving, which is the only
tell a replay attack leaves. Nine hand-crafted stability features
(autocorrelation, trend slope, and so on) ride alongside the
learned \gru{} state.

\textbf{A mixture-of-experts defender.} One set of weights asked to
recognise five different attack signatures ends up mediocre at the
hard ones. Splitting into six specialist experts with a learned
router fixes that. On Modena the MoE defender reaches anomaly F1
0.725 at AUROC 0.963, and this is the model self-play starts from.

\textbf{Adversarial self-play.} An Attacker GNN is co-trained
against the defender in an alternating Stackelberg loop: the
attacker makes sparse, budget-bounded, physics-respecting
perturbations, the defender is updated to catch them, then they
swap. Turning the attacker into its own mixture of four specialists,
with a diversity loss and load balancing, makes the population split
into two attack families on its own, with no class labels anywhere.
The defender that comes out of this is 5.8 points of F1 better than
the pretrained one on Modena, 23\% better on pressure MAE, and,
more to the point, better on two attack types it was never shown.

\section{Scope and Limitations}
\label{sec:scope}

The thesis focuses on sensor-level anomaly detection, not on
identifying the attacker's identity or intent. The threat model
assumes a white-box attacker who has read access to the network
topology and sensor feature space but cannot modify the model
weights. Physical attacks such as pipe tampering or pump sabotage
are outside scope.

The networks studied (Net1, Net3, Modena) are standard EPANET
benchmarks; they do not have publicly available ground-truth attack
logs, so all attacks are synthetically generated. This is common
practice in the literature but means that real-world deployment
would require re-evaluation on field data.

Three honest limitations of the results are documented in
Chapter~\ref{ch:conclusions}: replay attacks remain partially
undetected, self-play hurts the smallest network (Net1), and the
Attacker \moe{} uses only two of four experts at convergence.

\section{Thesis Structure}
\label{sec:structure}

Chapter~\ref{ch:background} reviews the necessary background: water
distribution network hydraulics, the cyber-attack threat model,
graph neural network architectures, temporal sequence models, and
adversarial machine learning. Chapter~\ref{ch:problem} formalises
the problem, defines the attack taxonomy, and describes the three
benchmark networks and the dataset generation procedure.
Chapter~\ref{ch:architecture} details the four stages of the
architecture progression. Chapter~\ref{ch:selfplay} presents the
Stackelberg self-play framework, the attacker model, the
auto-curriculum, and the Mixture-of-Attackers extension.
Chapter~\ref{ch:results} reports experimental results across all
three networks, with comparisons against classical baselines, an
ablation over \gnn{} backbone types, and the held-out novelty
experiments. Chapter~\ref{ch:conclusions} discusses limitations and
points to future work.
